%
% Documento: Resumo (Português)
%

\begin{center}
\imprimirtitulo
\end{center}

\begin{resumo}
O desenvolvimento de aplicações front-end modernas demanda ecossistemas que equilibrem agilidade de entrega, facilidade de manutenção e confiabilidade de código. Nesse cenário, o debate entre o uso de JavaScript (JS), com sua tipagem dinâmica tradicional, e de TypeScript (TS), com seu sistema de tipagem estática opcional, tornou-se central na engenharia de software. Este Trabalho de Conclusão de Curso investiga as percepções e os fatores determinantes na adoção de JavaScript e TypeScript, analisando comparativamente as diferenças entre desenvolvedores iniciantes e profissionais experientes. A pesquisa adotou uma abordagem mista (quantitativa e qualitativa) por meio de um estudo de levantamento (\textit{survey}) aplicado a 23 desenvolvedores, contemplando desde iniciantes até especialistas com vivência em projetos corporativos de grande escala. Os dados foram analisados através de estatística descritiva e categorização temática de respostas abertas. Os resultados revelam uma marcante divergência cognitiva condicionada pela maturidade profissional: enquanto iniciantes priorizam a flexibilidade e a prototipagem rápida do JavaScript (média de 4,29 na escala Likert), profissionais experientes manifestam ganhos expressivos de produtividade com o TypeScript (média de 4,22 contra 2,93 de iniciantes) e esmagadora preferência por sua adoção em novos projetos corporativos (4,56). Ambos os grupos convergem no reconhecimento de que a ausência de tipagem no JavaScript é propensa à geração de falhas em tempo de execução (média de 4,17) e que o TypeScript aprimora substancialmente a detecção precoce de erros (4,17 geral e 4,78 entre seniores), a manutenibilidade a longo prazo (4,44) e a padronização em equipes multidisciplinares (4,30). Constata-se que a decisão de adoção transcende aspectos puramente sintáticos, configurando-se como uma escolha arquitetural estratégica influenciada pela escala do projeto e pela experiência da equipe de desenvolvimento.

\vspace{\onelineskip}
\noindent
\textbf{Palavras-chave}: JavaScript. TypeScript. Desenvolvimento Front-End. Tipagem Estática. Manutenibilidade de Software. Engenharia de Software.
\end{resumo}
